\chapter{What the average keeps}

This chapter shows that the field is continuous and the molecules sit in it, that equation (1) is
a statement about averages of both over cells, that an average over enough volume and enough time
gives an entropic equivalence that is approximate and never exact, and that a solution which breaks
down takes away both the volume and the time. It rests on the second observation of the thought
experiment, Proposition 5 of the previous chapter, and the density and viscosity of water from
IAPWS R6-95 (2018)~\cite{src:IAPWS-R6-95-2018} and IAPWS R12-08~\cite{src:IAPWS-R12-08}.

\section{The observation}

Take a quintillion molecules of a fluid. Each one feels the field and every neighbor, and each one
moves on its own path. Equation (1) does not hold any of those paths. It holds a density, a
velocity and, once the temperature is put back, a temperature: a few numbers at each point,
standing for every molecule near that point.

\section{The field and the molecules}

The field is defined at every point of $\R^3$. The molecules sit at a finite number of points in it.
Nothing about the one rules out the other.

\textbf{Proposition 8.} Let $\Phi$ be a continuous field on $\R^3$ and let $x \neq y$ be any two
points, two molecules among them.
\begin{enumerate}
\item If $\Phi(x) \neq \Phi(y)$, then on the segment from $x$ to $y$ the field takes every value
between $\Phi(x)$ and $\Phi(y)$, and those values are uncountably many.
\item The map $g(t) = t/(1+t)$ carries $[0,\infty)$ one to one onto $[0,1)$ and keeps order. Every
finite value, however large, has a place in $[0,1)$, and none reaches $1$.
\end{enumerate}

\emph{Proof.} The function $s \mapsto \Phi(x + s(y-x))$ is continuous on $[0,1]$, and the
intermediate value theorem gives every value between its ends. An interval of real numbers with
two distinct ends is uncountable. For the second, $g'(t) = 1/(1+t)^2 > 0$, $g(0) = 0$, $g(t) < 1$,
$g(t) \to 1$, and $s \mapsto s/(1-s)$ is its inverse on $[0,1)$. \hfill\qedsymbol

The differences of the field between two molecules are uncountably many, and every one of them is
finite. Normalized to $[0,1)$ they are all kept and none is cut off. The previous chapter is about
this field, which is continuous at every scale. This chapter is about what equation (1) does with
it: it averages the field, and every molecule in it, over a cell.

The average is in the definition of the velocity. Stokes~\cite{src:Stokes-1845}, Article 1: ``we
may consider the mean velocity of the molecules ... apart
from the velocity due to the irregular motion. It is this regular velocity which I shall understand
by the velocity of a fluid at any point.''

\section{The average keeps a class and drops its members}

Take $N \ge 2$ molecules of mass $m$ in one cell, with velocities $v_1,\dots,v_N$. The cell values the
balance laws carry are the count $N$, the momentum $P = m\sum_i v_i$, and the kinetic energy
$K = \tfrac12 m\sum_i |v_i|^2$. The count below is for the kinetic part alone. Energy held in
rotation, vibration and the forces between molecules adds to it.

\textbf{Proposition 9.} For fixed $N$, $P$ and $K$ with $K > |P|^2/(2Nm)$, the velocities that give
those values form a sphere of dimension $3N-4$. With the positions free in the cell, the states of
the $N$ molecules that give the same cell values form a set of dimension $6N-4$.

\emph{Proof.} Write $v_i = P/(Nm) + w_i$ with $\sum_i w_i = 0$. The cross term in $K$ is
$\tfrac{1}{N}P\cdot\sum_i w_i = 0$, and
\[
K = \frac{|P|^2}{2Nm} + \frac{m}{2}\sum_i |w_i|^2 .
\]
The $w$ lie in a subspace of dimension $3N-3$ and on a sphere in it of positive radius, which has
dimension $3N-4$. The $3N$ coordinates of position add to it. \hfill\qedsymbol

For a quintillion molecules the cell keeps four numbers and drops a set of dimension
$6\times 10^{18} - 4$. Every state in that set gives equation (1) the same values. The set is the
equivalence class of the cell, and its size is what entropy measures: the Boltzmann entropy of the
cell values is a constant times the logarithm of the size of the set. The average keeps the class.
Which member of the class the molecules are in, which is every force on every molecule, is what it
drops.

\section{The equivalence is approximate at every finite count}

\textbf{Proposition 10.} Let $X_1,\dots,X_N$ be independent, each with mean $a$ and squared spread
$E[(X_i-a)^2] = s^2 > 0$, and let $\bar X = \tfrac1N\sum_i X_i$. Then $\bar X$ has mean $a$ and
squared spread $s^2/N$, which is positive for every finite $N$.

\emph{Proof.} The squared spread of a sum of independent terms is the sum of the squared spreads of the terms, since every
cross term $E[(X_i-a)(X_j-a)]$ with $i\neq j$ is zero, and scaling by $1/N$ scales it by $1/N^2$.
\hfill\qedsymbol

The average equals the value it stands for only in the limit of infinitely many molecules. For a
quintillion, its scatter is about $10^{-9}$ of the scatter of one molecule, and the continuum
works in a glass of water for that reason. It is approximate there, and very good, and never
exact.

\textbf{What Proposition 10 assumes.} It assumes independence. Molecules in a liquid are not
independent: each one is held by its neighbors. Correlation between them makes the scatter larger,
and the $1/\sqrt N$ rate is then the best case, not a bound. The same shape applies to an average
over time, with the time over which the motion stays correlated in the place of one molecule; that
form is stated and not proved here.

\section{Where there is not enough volume or time}

\textbf{Proposition 11.} In the class of Proposition 5, let
$\ell(t) = \|u(t)\|_{L^2}/\|\nabla u(t)\|_{L^2}$, a length over which the velocity changes. Then
$\ell$ has no positive lower bound on $[0,T^*)$, and neither does $\ell^2/\nu$, the time viscosity
takes to act across $\ell$.

\emph{Proof.} $\ell(t) \le \sup_s\|u(s)\|_{L^2}/\|\nabla u(t)\|_{L^2}$. The numerator is finite by
hypothesis, and Proposition 5 shows the denominator has no bound. \hfill\qedsymbol

Equation (1) carries no molecules, and Proposition 11 is a statement about equation (1) alone.
Read as an average of a real fluid, it asks for cells that hold enough molecules and for a time
over which to average that is long against the motion of one molecule and short against the motion
of the flow. Proposition 11 takes both away. The length of the flow goes to zero, and with it the
count of molecules in a cell of that length. The time of the flow goes to zero, and with it the
window in which an average over time can be taken.

For water at $20\,^\circ$C and atmospheric pressure, IAPWS R6-95 (2018) gives a density of
998.2072 kilograms per cubic meter, and with its specific gas constant and the Boltzmann constant
that is $3.337\times 10^{28}$ molecules per cubic meter. A cube of side one micrometer holds
$3.34\times 10^{10}$ of them, with a scatter under independence of $5.5\times 10^{-6}$. A cube of
side one nanometer holds $33.4$, with a scatter of $17$ percent. IAPWS R12-08 gives
$\nu = 1.0034\times 10^{-6}$ square meters per second, and $\ell^2/\nu$ is one picosecond at that
length. A solution of equation (1) that breaks down passes through
those lengths and those times on its way, and there the average it is a statement about is no
longer a good one.

\section{What this establishes and what it does not}

\textbf{Established.} Equation (1), read as a law of a real fluid, is a law of cell averages. The
average keeps an equivalence class of dimension $6N-4$ and drops which member of it the molecules
are in. Under independence, the average is approximate at every finite count and exact only in the
limit. In the class of the forced construction, the length and the time of the flow have no
positive lower bound. The cells and the windows the average needs do not stay large enough.

\textbf{Not established.} Nothing here says what the molecules do in a collapsing core, or that a
description that keeps them would not break down. The scatter figures rest on independence, which
a liquid does not have. Proposition 9 counts only the states of the molecules. The values the field
takes between them, uncountably many by Proposition 8, are dropped by the same average and are not
in that count.
